\section{Complexidade, truncagem e estabilidade}
\label{sec:complexity-dp}

\subsection{Contagem de operações}

Escreva
\[
 M=|\Values_{D,H}(K)|,\qquad S_D=2^{D+1}-1.
\]
As cotas abaixo usam o modelo de custo unitário para operações e comparações
no corpo ordenado dos custos. O tamanho dos operandos racionais será discutido
separadamente. O custo do terceiro componente da chave \eqref{eq:tie-key}, o
\emph{fluxo binário completo}, é contado em bits. Comparar duas
chaves que empatam nas duas primeiras componentes custa
$O(L)$, onde
\begin{equation}
 L=\max_R|\ser_{\rm rec}(R)|
 \le S_D\bigl(2\lfloor\log_2|\Library_H|\rfloor+2\bigr)
 \label{eq:stream-length}
\end{equation}
é o comprimento máximo de fluxo no catálogo. Todas as cotas desta seção contam
comparações de chave e construção/concatenação de fluxos como unitárias;
multiplicá-las por $L$ dá uma cota superior correspondente no modelo em que
comparar ou copiar um bit tem custo unitário.
Pela mesma razão, as cotas de memória contam \emph{entradas} da tabela, e cada
entrada que armazena um fluxo ocupa até $L$ bits. Uma implementação que
guarde apenas ponteiros de reconstrução e recomponha o fluxo sob demanda
troca esse fator de memória por trabalho adicional na comparação.

\begin{theorem}[complexidade em número de estados]
\label{thm:state-complexity}
Suponha os custos de folha pré-computados. Uma implementação densa do
\cref{alg:variable-tree} usa
\begin{equation}
 O(S_DM^3)
 \label{eq:dense-time}
\end{equation}
operações e
\begin{equation}
 O(S_DM^2)
 \label{eq:dense-memory}
\end{equation}
entradas de memória quando preserva todos os apontadores de reconstrução.
A versão com orçamento de até $N$ folhas usa
\begin{equation}
 O(S_DN^2M^3)
 \label{eq:budget-time}
\end{equation}
operações e $O(S_DNM^2)$ entradas. Se cada custo local
\eqref{eq:local-loss} for calculado diretamente sobre $P$ observações, a
inicialização densa acrescenta $O(S_DPM^2)$ operações.
\end{theorem}

\begin{proof}
Cada nó tem no máximo $M^2$ estados externos $(u,v)$. Em um nó interno, a
recorrência testa no máximo $M$ valores intermediários $w$ por estado, logo
$O(M^3)$ operações. Há no máximo $S_D$ nós. Armazenar um custo e um apontador
por estado dá \eqref{eq:dense-memory}.

Com o número de folhas no estado, há $N$ camadas de custo. Somando, para todos
$n\le N$, as decomposições positivas $n_L+n_R=n$, obtêm-se $O(N^2)$ pares de
orçamentos por tripla $(u,w,v)$. Isso prova \eqref{eq:budget-time} e a cota de
memória. A última afirmação vem de somar até $P$ termos em cada custo local.
\end{proof}

\begin{remark}[tabelas esparsas]
\label{rem:sparse}
As cotas são deliberadamente densas. O algoritmo distribuído armazena apenas
pares realizáveis e agrupa o filho direito por $w$. Seu trabalho em um nó é
proporcional ao tamanho da junção
\[
 \#\{((u,w),(w,v)):\text{os dois estados estão presentes}\},
\]
que pode ser muito menor que $M^3$. Essa melhora depende da instância e não
altera a cota de pior caso.

Convém dizer o quanto. Na instância principal da \cref{sec:numerical}, com
$(D,H)=(3,4)$, o alfabeto efetivo tem $M=475$ elementos, contra
$2S_Dq_H(q_H-1)=15\,180$ pela primeira cota de
\eqref{eq:value-alphabet-bound} e $36H^6=147\,456$ pela segunda; e a tabela da
raiz contém $3\,383$ estados realizáveis, contra $M^2=225\,625$ posições
densas, isto é, uma ocupação de $1{,}5\%$. Nas instâncias examinadas, a
estrutura esparsa reduz substancialmente o trabalho em relação à cota densa de
pior caso.
\end{remark}

\subsection{Modelo de custo e complexidade de bits}

Pelo \cref{lem:value-alphabet},
\[
 M\le36H^6,
\]
e portanto \eqref{eq:dense-time} fornece, de modo bruto,
\begin{equation}
 O(S_DH^{18})
 \label{eq:height-worst}
\end{equation}
operações de custo unitário. A expressão é polinomial no valor numérico de
$H$, cujo comprimento binário é $O(\log H)$, e
$S_D=2^{D+1}-1$. Trata-se, portanto, de uma cota pseudo-polinomial nos
orçamentos codificados em binário.

Para fechar também esse segundo nível, defina, para $r=p/q$ reduzido, $q>0$,
\begin{equation}
 \operatorname{bl}(r)
 :=1+\left\lceil\log_2(|p|+1)\right\rceil
    +\left\lceil\log_2(q+1)\right\rceil .
 \label{eq:rational-bit-length}
\end{equation}
Seja $B_{\mathcal Z}$ o maior comprimento \eqref{eq:rational-bit-length}
entre abscissas, respostas, pesos e penalidades racionais da instância. Denote
por $\mathsf M_{\mathbb Z}(b)$ uma cota para multiplicar inteiros de $b$ bits,
e recorde o limite $L$ de \eqref{eq:stream-length}.

\begin{theorem}[cota grosseira de complexidade de bits]
\label{thm:bit-complexity}
Todos os custos racionais que aparecem na recorrência possuem numeradores e
denominadores de comprimento no máximo
\begin{equation}
 B_{\rm cost}
 =O\!\left(
 P\bigl(B_{\mathcal Z}+D+|k_0|+\log(2H)\bigr)
 +B_{\mathcal Z}+D+\log(2L)
 \right).
 \label{eq:cost-bit-growth}
\end{equation}
Consequentemente, a versão densa sem orçamento de folhas pode ser executada em
\begin{equation}
 O\!\left(
 (S_DM^3+S_DPM^2)
 \bigl[\mathsf M_{\mathbb Z}(B_{\rm cost})
       \log(2B_{\rm cost})+L\bigr]
 \right)
 \label{eq:dense-bit-time}
\end{equation}
operações de bit e
\begin{equation}
 O\!\left(S_DM^2(B_{\rm cost}+L)\right)
 \label{eq:dense-bit-memory}
\end{equation}
bits de memória quando todos os estados e fluxos são armazenados. Na versão
com até $N$ folhas, substituem-se $S_DM^3$ e $S_DM^2$ por
$S_DN^2M^3$ e $S_DNM^2$, respectivamente. Com multiplicação escolar,
$\mathsf M_{\mathbb Z}(b)=O(b^2)$.
\end{theorem}

\begin{proof}
Um átomo admissível tem $\operatorname{bl}(y_0),
\operatorname{bl}(s)=O(\log(2H))$. Em uma célula descendente de profundidade
no máximo $D$, escrever o centro ou um extremo como racional expande o fator
diádico em no máximo $D+|k_0|+O(1)$ bits. Portanto a avaliação afim em uma
abscissa da amostra tem comprimento
\[
 O(B_{\mathcal Z}+D+|k_0|+\log(2H)).
\]
Subtrair a resposta, elevar ao quadrado e multiplicar pelo peso apenas altera
a constante dessa cota. Na soma ingênua de $P$ frações, pode-se tomar como
denominador o produto dos denominadores; logo numerador e denominador têm
comprimento no máximo $P$ vezes o maior comprimento de uma parcela. Os termos
de penalização acrescentam $O(B_{\mathcal Z}+D+\log(2L))$ bits. Isso prova
\eqref{eq:cost-bit-growth}; cancelamentos por máximo divisor comum só podem
reduzir os operandos.

Uma soma, comparação ou redução de frações com operandos de comprimento
$B_{\rm cost}$ custa
$O(\mathsf M_{\mathbb Z}(B_{\rm cost})\log(2B_{\rm cost}))$ operações de bit
por algoritmos inteiros usuais. Quando custos e número de folhas empatam, a
comparação shortlex lê no máximo $L$ bits. Multiplicar esse custo pelo número
de transições e inicializações do \cref{thm:state-complexity} dá
\eqref{eq:dense-bit-time}. Cada estado guarda um custo, um fluxo ou apontador
de reconstrução e índices de $O(\log(S_DM))$ bits, absorvidos pela cota exibida;
isso prova \eqref{eq:dense-bit-memory}. O argumento com a coordenada de
orçamento é idêntico.
\end{proof}

A cota é propositalmente conservadora: ela não explora cancelamentos nem
denominadores compartilhados. Ainda assim, fecha o modelo de Turing e torna
explícito o ponto essencial. Como $M\le36H^6$ e $S_D=2^{D+1}-1$, a
dependência continua pseudo-polinomial nos valores numéricos de $H$ e $2^D$,
mesmo após contabilizar o crescimento dos racionais.

\subsection{Valores sob truncagem}

Seja
\[
 \Representations_{\infty}^{\Library}(K)
 =\bigcup_{D,H}\Representations_{D,H}^{\Library}(K)
\]
o conjunto de todas as representações finitas enraizadas em $K$.

\begin{proposition}[convergência dos mínimos penalizados]
\label{prop:truncation-convergence}
Os valores $C_{D,H}^{\star}(K)$ de \eqref{eq:variable-global} são não
crescentes separadamente em $D$ e $H$, e
\begin{equation}
 \inf_{D,H}C_{D,H}^{\star}(K)
 =\inf_{R\in\Representations_{\infty}^{\Library}(K)}
   \cost_{\mathcal Z}(R).
 \label{eq:penalized-limit}
\end{equation}
Ao longo de qualquer sequência componente a componente não decrescente e
cofinal em $(D,H)$, os valores convergem ao lado direito.

Se o ínfimo em \eqref{eq:penalized-limit} é atingido, o minimizador global de
menor chave \eqref{eq:tie-key} pertence a algum catálogo finito; a partir de
uma truncagem que o contenha, o seletor canônico estabiliza.
\end{proposition}

\begin{proof}
Aumentar $D$ ou $H$ preserva toda representação já admissível, logo os
mínimos não crescem. A união dos domínios dos mínimos é exatamente
$\Representations_{\infty}^{\Library}(K)$, o que prova
\eqref{eq:penalized-limit}. Uma sequência cofinal eventualmente contém cada
representação finita; a monotonicidade conclui a convergência.

Se há minimizadores globais, seus comprimentos de fluxo formam um subconjunto
não vazio de $\Nzero$ e possuem mínimo. Entre as finitíssimas palavras desse
comprimento existe uma menor lexicograficamente. A representação
correspondente tem profundidade e altura finitas. Depois que ela entra no
catálogo, nenhum candidato posterior pode melhorar o custo ou a chave.
\end{proof}

\subsection{Um limite de estabilidade}

\begin{proposition}[um seletor exato pode ser descontínuo]
\label{prop:instability}
Considere, em $K=[0,1]$, o subcatálogo formado pelos dois átomos
\[
 g_+(x)=x,\qquad g_-(x)=-x,
\]
ambos de altura $2$. Para uma única observação $(x,z)=(1,z)$, custo igual à
perda quadrática e qualquer desempate determinístico, o mapa que envia $z$ ao
átomo selecionado é descontínuo em $z=0$ na norma uniforme.
\end{proposition}

\begin{proof}
Os custos são
\[
 \cost_z(g_+)=(z-1)^2,\qquad
 \cost_z(g_-)=(z+1)^2.
\]
Logo $g_+$ é o único minimizador para $z>0$, $g_-$ é o único minimizador para
$z<0$ e há empate em $z=0$. Qualquer regra escolhe um deles no empate, mas a
sequência que se aproxima por dados do lado oposto seleciona sempre o outro.
Como
\[
 \|g_+-g_-\|_{\infty,[0,1]}=2,
\]
não há continuidade.
\end{proof}

O exemplo separa exatidão de continuidade do seletor. Hipóteses de margem
entre custos produzem regiões de seleção localmente constante; nos empates,
uma saída a valores conjuntos preserva semicontinuidade. Essa distinção entre
existência, estabilidade e não unicidade também aparece na análise de problemas
de treinamento por perda quadrática em esquemas não lineares
\citep{christof2021}.
